\documentclass[tikz,border=8pt]{standalone}
\usepackage{ctex,amsmath,amssymb}
\usetikzlibrary{arrows.meta,positioning,calc,fit,backgrounds}
% Shared visual language for LFS architecture diagrams.
\RequirePackage{../../mymath}

\definecolor{LFSBlue}{HTML}{2563EB}
\definecolor{LFSBlueSoft}{HTML}{DBEAFE}
\definecolor{LFSGreen}{HTML}{15803D}
\definecolor{LFSGreenSoft}{HTML}{DCFCE7}
\definecolor{LFSOrange}{HTML}{C2410C}
\definecolor{LFSOrangeSoft}{HTML}{FFEDD5}
\definecolor{LFSRed}{HTML}{B91C1C}
\definecolor{LFSRedSoft}{HTML}{FEE2E2}
\definecolor{LFSPurple}{HTML}{7E22CE}
\definecolor{LFSPurpleSoft}{HTML}{F3E8FF}
\definecolor{LFSGray}{HTML}{475569}
\definecolor{LFSGraySoft}{HTML}{F1F5F9}

\tikzset{
  lfs picture/.style={
    font=\small,
    node distance=8mm and 12mm,
    >=Latex,
    every path/.style={line cap=round, line join=round}
  },
  block/.style={
    draw=LFSBlue,
    fill=LFSBlueSoft,
    rounded corners=2pt,
    line width=0.8pt,
    align=center,
    inner xsep=7pt,
    inner ysep=5pt,
    minimum height=10mm
  },
  compute/.style={
    block,
    draw=LFSPurple,
    fill=LFSPurpleSoft
  },
  storage/.style={
    block,
    draw=LFSGreen,
    fill=LFSGreenSoft
  },
  interface/.style={
    block,
    draw=LFSOrange,
    fill=LFSOrangeSoft
  },
  risk/.style={
    block,
    draw=LFSRed,
    fill=LFSRedSoft
  },
  neutral/.style={
    block,
    draw=LFSGray,
    fill=LFSGraySoft
  },
  decision/.style={
    diamond,
    draw=LFSOrange,
    fill=LFSOrangeSoft,
    line width=0.8pt,
    align=center,
    inner sep=2pt,
    aspect=2
  },
  group/.style={
    draw=LFSGray,
    rounded corners=3pt,
    dashed,
    line width=0.7pt,
    inner sep=6pt
  },
  arrow/.style={
    -{Latex[length=2.4mm,width=1.6mm]},
    draw=LFSGray,
    line width=0.9pt
  },
  data arrow/.style={
    arrow,
    draw=LFSBlue
  },
  control arrow/.style={
    arrow,
    draw=LFSOrange,
    dashed
  },
  residual/.style={
    arrow,
    draw=LFSGreen,
    rounded corners=4pt
  },
  label text/.style={
    font=\scriptsize,
    text=LFSGray,
    fill=white,
    inner sep=1.5pt
  },
  title/.style={
    font=\bfseries,
    text=LFSGray,
    align=center
  }
}


\begin{document}
\begin{tikzpicture}[my picture, node distance=5mm and 6.5mm]

% 顶部左侧：结构演变与稀疏专家
\node[block, text width=68mm, minimum height=13mm] (moe_arch) {
  \textbf{混合专家架构（MoE）}\\[0.8mm]
  \scriptsize 稀疏门控路由机制、Top-$k$ 专家激活\\
  负载均衡辅助损失与专家并行（第27章）
};
\node[neutral, text width=68mm, minimum height=13mm, below=4mm of moe_arch] (bert_arch) {
  \textbf{双向掩码语言模型（BERT）}\\[0.8mm]
  \scriptsize 双向全向上下文编码、掩码语言建模（MLM）\\
  判别式与生成式表征对比（第25章）
};
\begin{scope}[on background layer]
  \node[group, fit=(moe_arch)(bert_arch), label={[font=\footnotesize\bfseries, text=MyBlue]above:一、结构拓展与稀疏计算}] (g_sparse) {};
\end{scope}

% 顶部右侧：长上下文与开源演进
\node[storage, text width=68mm, minimum height=13mm, right=8mm of moe_arch] (ctx_arch) {
  \textbf{长上下文扩展与外推}\\[0.8mm]
  \scriptsize 旋转位置编码插值（YaRN、外推缩放）\\
  超长序列记忆压缩与召回评估（第28章）
};
\node[compute, text width=68mm, minimum height=13mm, below=4mm of ctx_arch] (open_arch) {
  \textbf{开源主流架构演进谱系}\\[0.8mm]
  \scriptsize LLaMA、Mistral、Qwen、DeepSeek 设计演进\\
  归一化、激活函数与注意力变体权衡（第29章）
};
\begin{scope}[on background layer]
  \node[group, fit=(ctx_arch)(open_arch), label={[font=\footnotesize\bfseries, text=MyPurple]above:二、容量突破与开源架构谱系}] (g_open) {};
\end{scope}

% 底部：多模态理解与扩散生成
\node[storage, text width=45mm, minimum height=14mm, below=11mm of bert_arch.south west, anchor=north west] (vit_arch) {
  \textbf{视觉 Transformer（ViT）}\\[0.8mm]
  \scriptsize 图像二维分块（Patch Embedding）\\
  无偏置自注意力机制（第81章）
};
\node[interface, text width=45mm, minimum height=14mm, right=5mm of vit_arch] (mllm_arch) {
  \textbf{多模态对齐与融合（MLLM）}\\[0.8mm]
  \scriptsize 跨模态投影器（Projector）\\
  图文交织自回归统一建模（第82章）
};
\node[compute, text width=45mm, minimum height=14mm, right=5mm of mllm_arch] (diff_arch) {
  \textbf{扩散概率与 DiT 生成}\\[0.8mm]
  \scriptsize 正向加噪与反向去噪 SDE\\
  DiT 条件视觉生成架构（第83、84章）
};
\begin{scope}[on background layer]
  \node[group, fit=(vit_arch)(mllm_arch)(diff_arch), label={[font=\footnotesize\bfseries, text=MyOrange]above:三、多模态感知与跨模态生成演进}] (g_multi) {};
\end{scope}

% 连接线
\draw[data arrow] (bert_arch.north) -- (moe_arch.south);
\draw[data arrow] (ctx_arch.south) -- (open_arch.north);
\draw[data arrow] (vit_arch.east) -- (mllm_arch.west);
\draw[data arrow] (mllm_arch.east) -- (diff_arch.west);

\draw[data arrow] ($(bert_arch.south)+(15mm,0)$) -- ($(vit_arch.north)+(15mm,0)$);
\draw[data arrow] ($(open_arch.south)-(15mm,0)$) -- ($(diff_arch.north)-(15mm,0)$);

% 底部总结栏
\node[neutral, text width=146mm, minimum height=7mm, below=6mm of mllm_arch.south] (bot_summary) {
  \scriptsize \textbf{第五篇拓展主线：} 双向与稀疏结构创新 $\longleftrightarrow$ 百万长上下文突破 $\longleftrightarrow$ 开源生态架构融合 $\longleftrightarrow$ 视觉理解与扩散生成多模态拓展
};

\end{tikzpicture}
\end{document}
